\documentclass[11pt]{article}

\usepackage[margin=1in]{geometry}
\usepackage[T1]{fontenc}
\usepackage[utf8]{inputenc}
\usepackage{amsmath}
\usepackage{amssymb}
\usepackage{amsthm}
\usepackage{hyperref}
\usepackage{microtype}

\hypersetup{
  colorlinks=true,
  linkcolor=blue,
  citecolor=blue,
  urlcolor=blue,
  pdftitle={Reducing Floats mod p},
  pdfauthor={Benoit Jacob <benoit.jacob@amd.com>}
}

\newcommand{\Z}{\mathbb{Z}}
\newcommand{\F}{\mathbb{F}}
\newcommand{\Q}{\mathbb{Q}}
\newcommand{\R}{\mathbb{R}}
\newcommand{\ord}{\operatorname{ord}}
\newcommand{\lcm}{\operatorname{lcm}}
\newcommand{\im}{\operatorname{im}}
\newcommand{\Fp}{\F_p}
\newcommand{\FieldFast}{\texttt{FieldFast}}
\newcommand{\FieldMC}{\texttt{FieldWith\allowbreak{}Mul\allowbreak{}Casts}}

\theoremstyle{plain}
\newtheorem{constraint}{Constraint}
\newtheorem{proposition}{Proposition}

\title{Reducing Floats mod \(p\)\\
       \large Algebraic FPSan as a residue semantics}
\author{Benoit Jacob\\\texttt{benoit.jacob@amd.com}}
\date{}

\begin{document}
\maketitle

\begin{abstract}
FPSan is a dynamic testing technique for numerical programs: replace each
floating-point value by a same-width fingerprint, run the program in a finite
target ring, and compare the resulting fingerprints.  The original
Triton FPSan introduced this idea using a tuned bijection from source-format bit
patterns into \(\Z/2^w\Z\), where \(w\) is the source width
\cite{triton-fpsan-source,goucher-fpsan-blog}.  That scramble is useful, but it
is not a homomorphism from the represented numeric values, so identities such as
\(2+2=4\) or \(x+x=2x\) need not hold.

This paper studies algebraic FPSan: a homomorphic alternative to Triton's FPSan
scramble on finite binary floating-point values.  The observation is simple:
every finite binary float is a dyadic rational, hence an element of
\(\Z[1/2]\).  For every odd prime \(p\), reduction modulo \(p\) gives a ring
homomorphism
\[
  \Z[1/2] \longrightarrow \F_p.
\]
This gives a finite-field residue semantics for FPSan that preserves exact
rational identities while keeping the same-width fingerprint invariant.  We then
study the constraints on the prime \(p\), multiplicative casts between widths,
composite residue rings for selected exponential, logarithmic, and trigonometric
laws, multi-prime recovery by the Chinese remainder theorem (CRT), and
extensions to algebraic-number formats.
\end{abstract}

\section{Introduction}

Suppose a numerical kernel is rewritten for performance.  A reduction may be
reassociated, an expression may be factored, an FMA may replace a multiply and an
add, or a matrix intrinsic may change the accumulation order.  Comparing
floating-point results bit-for-bit is too strict: a value-preserving algebraic
rewrite can change rounding.  Tolerance tests, on the other hand, are often too
weak and too problem-specific.

Triton's FPSan introduced a useful third option
\cite{triton-fpsan-source}.  Instead of asking whether two rounded
floating-point executions are bitwise equal, FPSan asks whether the two
computations have the same fingerprint in a chosen target ring.  Each
source value is replaced by a fingerprint, and each source operation is replaced
by a deterministic operation on fingerprints.  The resulting semantics are
designed to be compact and invariant under formal laws such as associativity and
distributivity.  They are not a reference implementation of floating-point
rounding.

The invariant throughout this document is width preservation.  If the original
floating-point type has \(w\) bits, then every FPSan semantics discussed here
stores its fingerprint in exactly \(w\) bits.  The central design choice is the
target ring together with the leaf map.  One chooses a target ring encodable in
the same bit width, possibly with distinguished sentinel codes adjoined to the
finite ring, and one chooses how each source floating-point value is sent into
that target ring.

Triton's concrete semantics instantiate this design by choosing the target ring
\(\Z/2^w\Z\) and a tuned bijection, not a ring homomorphism,
\[
  E:\{\text{source-format bit patterns of width }w\}\longrightarrow \Z/2^w\Z.
\]
The construction itself is mostly indifferent to the
details of a particular floating-point standard: what matters is that the
sanitizer receives fixed-width bit patterns.  Writing \(0\) and \(1\) for the
bit patterns of those values, and writing \(-x\) for source-format sign
negation of a bit pattern \(x\), Triton's tuning gives \(E(0)=0\), \(E(1)=1\),
and \(E(-x)=-E(x)\) in the target ring.  After applying \(E\), arithmetic is
performed with ring operations in \(\Z/2^w\Z\), with special handlers for
operations such as division, exponentials, and trigonometry.  The Triton
author's blog discusses the algebraic intentions of this construction
\cite{goucher-fpsan-blog}.

Because \(E\) is a bijection, one can always define an artificial operation on
bit patterns by transporting the target-ring operation back through \(E\):
\[
  x \boxplus_E y = E^{-1}(E(x)+E(y)).
\]
With respect to \(\boxplus_E\), the map \(E\) is a homomorphism by construction.
That statement is true but not the statement a compiler engineer usually wants:
\(\boxplus_E\) is not floating-point addition, nor exact rational addition of
the represented values.  It is the arithmetic manufactured by the embedding.

The visible consequence is immediate.  If \(u\) is the source-format bit pattern
of \(2.0\), then typically
\[
  E(u) + E(u) \ne E(\text{bit pattern of }4.0).
\]
So Triton FPSan sees reassociation and distributivity, but it does not see
ordinary value facts such as \(2+2=4\), \(x+x=2x\), or \(x/x=1\) for every
nonzero \(x\).

The contribution of algebraic FPSan is that the opposite leaf-map choice is
possible.  We can keep compact same-width fingerprints, but choose the leaf map
to be a ring homomorphism on exact finite values.  The price is that a single
same-width reduction is no longer injective.  The benefit is that its collisions
are ordinary modular collisions in a characteristic-zero calculation reduced
modulo a prime, not artifacts of an unrelated transported arithmetic.  The next
section gives the algebraic observation that makes this possible.

In the \texttt{hip-fpsan} C++ API, these choices appear as the
\texttt{Semantics} enum.  Its enumerators are
\[
\begin{gathered}
\texttt{Native},\quad
\texttt{Triton},\\
\texttt{Field},\quad
\texttt{Field2},\quad
\texttt{FieldFast},\quad
\texttt{FieldFast2},\quad
\texttt{FieldWithMulCasts},\quad
\texttt{FieldWithMulCasts2},\\
\texttt{SophieGermainRing},\quad
\texttt{SophieGermainRing2},\quad
\texttt{PythagoreanRing},\quad
\texttt{PythagoreanRing2}.
\end{gathered}
\]
\texttt{Native} is ordinary floating-point arithmetic.  \texttt{Triton} is the
Triton FPSan semantics discussed above in this introduction.  All remaining
enum values are collectively called algebraic FPSan.  The Field-family semantics
\texttt{Field}, \texttt{FieldFast}, and \texttt{FieldWithMulCasts}, together
with their \(2\)-suffixed variants, are discussed in Sections~2--6.
\texttt{SophieGermainRing} and \texttt{SophieGermainRing2} are discussed in
Section~7, and \texttt{PythagoreanRing} and \texttt{PythagoreanRing2} in
Section~8.  A \(2\)-suffix means a second choice of primes or moduli under the
same design constraints: it rerolls the collision set without changing the
intended algebraic contract.

\section{Finite binary floats live in \texorpdfstring{\(\Z[1/2]\)}{Z[1/2]}}

A finite binary floating-point value can be written as
\[
  x = 2^e M,
\]
where \(M\in\Z\) is a signed integer significand and \(e\in\Z\).  Normal and
subnormal values both have this form, after absorbing the sign into \(M\).  Thus
every finite value lies in
\[
  \Z[1/2] = \left\{\,\frac{a}{2^k}:a\in\Z,\ k\ge 0\,\right\}.
\]
This is the localization of \(\Z\) obtained by inverting \(2\).

The maximal ideals of \(\Z[1/2]\) are exactly the ideals generated by odd primes.
Equivalently, an odd prime \(p\) is precisely a prime modulo which \(2\) remains
invertible.  The resulting residue homomorphism
\[
  \Z[1/2]\longrightarrow\F_p
\]
is the leaf map used by the Field-family semantics: a finite source value is
sent to its residue in the target field.  Concretely, this homomorphism is
determined by the image of \(1/2\), which is taken to be the inverse of \(2\) in
\(\F_p\).  This is the only algebraic fact about finite binary floats that the
Field-family semantics need.

This is the mathematical core of the Field-family semantics.  The fingerprint of
a finite value is the residue of its exact dyadic value.  Addition, subtraction,
and multiplication are the field operations.  The faithful Field variants use
field inversion for division; \FieldFast{} keeps the same finite-value
fingerprints and the same additive/multiplicative arithmetic, but treats
nontrivial division and roots as tagged operations for speed.  The cheap
division cases \(0/x=0\) and \(x/1=x\) are still preserved.

\section{Compactness, sentinels, and non-injectivity}

By width preservation, a source format of width \(w\) gives only \(2^w\)
fingerprint codes.  Let \(N\) be the finite residue modulus used by an algebraic
semantics; for example, in the \texttt{Semantics::Field} case with target
\(\F_p\), \(N=p\), while in the composite ring cases \(N=pd\).  The finite
residues use the codes \(0,\ldots,N-1\), and two further codes are reserved for
non-finite values:
\[
  N \quad\text{for infinity},\qquad N+1 \quad\text{for NaN}.
\]
Thus the finite modulus must satisfy \(N\le 2^w-2\).

This compactness has a price: the algebraic map is not injective.  It cannot be,
because many finite floats must share a finite set of residues.  This is unlike
Triton's bit-pattern scramble, which is a bijection before arithmetic starts.

The following fact is what makes that compromise much less dangerous at the most
important point.

\begin{proposition}[Zero is not a finite collision]
\label{prop:zero}
Let \(N\) be odd, and suppose every finite value of the source format can be
written as \(2^e M\) with \(|M|<N\).  Then its residue modulo \(N\) is zero if
and only if the represented value is zero.
\end{proposition}

\begin{proof}
Since \(N\) is odd, \(2\) is a unit modulo \(N\).  Hence
\[
  2^e M \equiv 0 \pmod N
  \quad\Longleftrightarrow\quad
  M\equiv 0 \pmod N.
\]
The hypothesis \(|M|<N\) makes the latter congruence equivalent to \(M=0\).
That is exactly the zero value; the two IEEE-754 signed zeros are deliberately
identified by the algebraic semantics.
\end{proof}

Nonzero collisions remain ordinary finite-modulus fingerprint collisions.  If
two distinct exact dyadic expressions on the same inputs compare equal, then
after clearing powers of \(2\), their difference has numerator divisible by the
chosen modulus.  The suffix-2 variants exist so that a suspected collision can
be checked against a different set of primes or composite moduli:
\texttt{Field2}, \texttt{FieldFast2}, \texttt{SophieGermainRing2}, and
\texttt{PythagoreanRing2} usually have different nonzero collision sets.  The
8-bit composite variants deliberately share the only 2-positive modulus that
fits their construction.

\section{Several primes and characteristic-zero recovery}

The non-injectivity of a single reduction should not be confused with loss of
the characteristic-zero arithmetic.  If \(p_1,\ldots,p_r\) are distinct odd
primes and \(P=\prod_i p_i\), then the combined map can be displayed as
\[
  \Z[1/2]\longrightarrow \prod_i \F_{p_i}
  \xrightarrow{\mathrm{CRT}} \Z/P\Z.
\]
The second arrow is Chinese remainder theorem (CRT) reconstruction, identifying
the product of residue fields with reduction modulo \(P\), where \(2\) is still
invertible.  Thus several algebraic FPSan runs can be interpreted as a single
wider residue computation.  If an exact dyadic result is known to be \(a/2^k\) with
\(|a|<P/2\), then the residues determine \(a\) uniquely by the symmetric CRT
representative, and hence determine the dyadic value itself.

This is a fundamental distinction from Triton's FPSan.  Triton's leaf map is
invertible, but after that the arithmetic is the arithmetic of the scrambled
ring \(\Z/2^w\Z\).  There is no compatible family of reductions whose CRT lift
is the original characteristic-zero calculation.  Algebraic FPSan goes the other
way: one prime is not injective, but different primes are all shadows of the
same calculation in \(\Z[1/2]\).

As a concrete bound, consider OCP FP8 E5M2 \cite{ocp-ofp8}.
It has exponent bias \(15\) and two mantissa bits.  Every finite value is
\(A/2^{16}\), and the maximum finite magnitude is
\[
  1.75\cdot 2^{15}=7\cdot 2^{13},
\]
so
\[
  |A|\le 7\cdot 2^{29}=3\,758\,096\,384.
\]
To reconstruct an arbitrary finite FP8 E5M2 value by signed CRT it is enough
to have \(P>2\cdot 7\cdot 2^{29}=7\,516\,192\,768\).  If each residue is still
required to fit in an 8-bit fingerprint with two sentinels, then \(p\le 254\).
The four largest possible primes give
\[
  251\cdot 241\cdot 239\cdot 233
  =3\,368\,562\,317,
\]
which is not enough; no four-prime set can do better.  Five primes suffice, for
example
\[
  251\cdot 241\cdot 239\cdot 233\cdot 229
  =771\,400\,770\,593.
\]
Thus five 8-bit prime reductions are cardinality-minimal for reconstructing any
finite rounded FP8 E5M2 value.  If one also insists on the Field-family
root congruence \(p\equiv 11\pmod {12}\), five primes still suffice, for example
\[
  251,\,239,\,227,\,191,\,179.
\]

The present library does not implement multi-prime reconstruction; its goal is
compact same-width fingerprinting.  The point is conceptual: algebraic FPSan
collisions are ordinary modular collisions, and ordinary CRT arithmetic can, in
principle, remove them within any known magnitude bound.  Rounding is a separate
semantic choice: the algebraic fingerprints track the exact dyadic expression
between whatever rounded values the instrumented program presents as leaves.

\section{Prime choices for the Field-family semantics}

The statement ``reduce modulo an odd prime'' leaves many possible primes.  The
implementation chooses primes satisfying several constraints simultaneously.
These constraints are motivated by different public semantics: some are needed
by the faithful \texttt{Field} root operations, some by
\FieldMC{} casts, and some by all Field-family modes.  The prime tables are
nevertheless shared by \texttt{Field}, \FieldFast{}, and \FieldMC{} (and by their
suffix-2 variants) so users can switch among these policies without changing the
finite fingerprints of inputs.

\subsection{Size and sentinel constraints}

The prime should be close to \(2^w\), because a larger field gives a lower
collision probability.  It must still leave room for infinity and NaN:
\[
  p\le 2^w-2.
\]
It must also be larger than the maximum possible \(|M|\) for the format's signed
significand \(M\), so that Proposition~\ref{prop:zero} applies and zero is hit
only by true zero.

\begin{constraint}[size, sentinels, and zero]
\label{con:size}
For a source format of width \(w\), the Field prime \(p\) must satisfy
\[
  M_{\max}(w)<p\le 2^w-2,
\]
where \(M_{\max}(w)\) is the largest possible \(|M|\) in the representation
\(x=2^eM\) for that format.
\end{constraint}

\subsection{Algebraic roots}

The faithful algebraic modes try to treat square root and cube root as algebraic
power maps when the chosen modulus allows it, but the exact contract differs by
semantics.  In the Field family, Constraints~\ref{con:sqrt}
and~\ref{con:cbrt} are used by \texttt{Field} and \FieldMC{}.  \FieldFast{}
shares the same primes for compatibility, but deliberately represents roots by
tagged fingerprints.

For \texttt{Semantics::Field}, the modulus is a prime \(p\), so the unit group
has order \(p-1\).  The square-root power map uses the following condition.

\begin{constraint}[square-root power map]
\label{con:sqrt}
\[
  p\equiv 3\pmod 4.
\]
This makes \(2\) invertible in the odd part of the unit-group exponent and gives
the multiplicative square-root choice used below.
\end{constraint}

Under this condition,

\[
  S(x)=x^{(p+1)/4}
\]
is a multiplicative square-root choice:
\[
  S(xy)=S(x)S(y).
\]
It is a true inverse to squaring on the square residues and gives a consistent
choice on the nonsquares.  No multiplicative square root can be a two-sided
inverse on all of \(\F_p^\times\), since squaring is two-to-one.

For cube roots, the power map \(x\mapsto x^r\) is a perfect cube root exactly
when \(3r\equiv 1\pmod {p-1}\).  That requires \(3\nmid p-1\), i.e.
\(p\equiv 2\pmod 3\).

\begin{constraint}[perfect cube-root power map]
\label{con:cbrt}
\[
  p\equiv 2\pmod 3.
\]
Equivalently, \(3\) is invertible modulo \(p-1\), so a power map can be an
inverse to cubing on \(\F_p^\times\).
\end{constraint}

Together, Constraints~\ref{con:sqrt} and~\ref{con:cbrt} give
\[
  p\equiv 11\pmod {12}.
\]
That congruence is the root-related constraint on the Field primes.  The Field
families do not additionally require \(2\) to be a quadratic residue; the prime
choices also serve density, same-width sentinel, and cast-tower constraints, so
the qr-order sign of \(2\) is mixed across widths and variants.  The composite
families below have another factor available and use it to make \(2\) positive
in the selected qr-order factor.

The composite algebraic modes use the same power-map idea with the exponent of
the unit group of \(\Z/N\Z\).  The Sophie Germain choices keep \(3\) invertible
in that exponent, so \texttt{cbrt} remains a perfect power-map cube root.  The
Pythagorean choices cannot keep that property while also making trigonometric
rotations available; there \texttt{cbrt} is represented as a tag.
Square root remains a multiplicative power map in the faithful algebraic modes,
although the round trip \(\sqrt{x}^2=x\) holds only on the part of the ring
selected by that power-map choice.

\subsection{Multiplicative casts}

The optional multiplicative-cast policy is the only cast policy here that tries
to preserve algebraic structure.  It is exposed as
\[
  \texttt{Semantics::FieldWithMulCasts}.
\]
The ordinary \texttt{Semantics::Field} variants and \FieldFast{} use the same
primes and finite-value fingerprints, but cheap deterministic casts.  Triton's
casts are signed integer resizes of the scrambled fingerprint, and the composite
algebraic semantics use a deterministic finite-residue convention.

A cast between different Field widths cannot be additive, hence cannot be a
ring homomorphism.  Indeed, if \(p\ne q\), every additive-group homomorphism
\((\F_p,+)\to(\F_q,+)\) is trivial.  The only nontrivial algebraic structure
available across different characteristics is therefore multiplicative: the
groups of nonzero residues.

The multiplicative construction below applies only to finite nonzero
fingerprints.  Zero is mapped to zero, and the infinity/NaN sentinels map to the
corresponding sentinels.  The \texttt{FP32}\(\leftrightarrow\)\texttt{FP64} edge
uses Pohlig--Hellman discrete logs because a linear scan over a group of size
about \(2^{32}\) would be too expensive; the group orders are chosen to make this
feasible.

\begin{constraint}[computable cast logarithms]
\label{con:cast-log}
The Field group orders used for cast logarithms, especially
\(m_{32}=p_{32}-1\), must be smooth enough for the Pohlig--Hellman path used by
the implementation.
\end{constraint}

The multiplicative group \(\F_p^\times\) is cyclic.  The multiplicative-cast
tower consists of the widths \(4,8,16,32,64\); the 6-bit format uses a
standalone prime and is not part of this tower.  Write \(p_i\) for the Field
prime at a tower width and \(m_i=p_i-1\).  For each format class \(\alpha\) at
that width, choose a primitive root \(g_{i,\alpha}\) of
\(\F_{p_i}^\times\).  There is usually only one such class; the important
exceptions are the two FP8 families, E4M3 and E5M2, and the two 16-bit
families, FP16 and BF16.  The implementation uses \(g_i\) for one class and
\(g_i^3\) for the other.  Since the Field primes are \(11\pmod {12}\), \(3\) is
coprime to \(m_i\), so \(g_i^3\) is again a primitive root.  Let
\[
  L_{i,\alpha}:\F_{p_i}^{\times}\to\Z/m_i\Z
\]
be discrete logarithm to the base \(g_{i,\alpha}\).  In these coordinates,
multiplication in \(\F_{p_i}^{\times}\) becomes addition in \(\Z/m_i\Z\).

The chosen primes satisfy
\[
  m_4 \mid m_8 \mid m_{16} \mid m_{32} \mid m_{64}.
\]
Moreover, the factors introduced at successive steps are coprime to the factors
already present.  Equivalently, whenever \(m_i\mid m_j\), the cofactor
\(m_j/m_i\) is coprime to \(m_i\), and stepwise cofactors between \(i\) and
\(j\) are pairwise coprime.

\begin{constraint}[multiplicative cast tower]
\label{con:cast-tower}
The Field primes in the cast tower must satisfy
\[
  m_4 \mid m_8 \mid m_{16} \mid m_{32} \mid m_{64},
  \qquad m_i=p_i-1,
\]
and whenever \(i<j\), the cofactor \(m_j/m_i\) must be coprime to \(m_i\); in
particular, the stepwise cofactors from \(i\) to \(j\) are pairwise coprime.
\end{constraint}

Here and below, \(i<j\) means that the \(i\)-bit tower format is narrower than
the \(j\)-bit tower format.  Put \(c_{ij}=m_j/m_i\).  The coprimality condition
in Constraint~\ref{con:cast-tower} gives a Chinese-remainder section
\[
  \sigma_{ij}:\Z/m_i\Z\to\Z/m_j\Z
\]
characterized by
\[
  \sigma_{ij}(a)\equiv a\pmod {m_i},\qquad
  \sigma_{ij}(a)\equiv 0\pmod {c_{ij}}.
\]
The widening (``up'') cast \(U_{i\alpha,j\beta}\) and narrowing (``down'') cast
\(D_{j\beta,i\alpha}\) are the multiplicative maps defined in logarithm
coordinates by
\[
  L_{j,\beta}(U_{i\alpha,j\beta}(x))=\sigma_{ij}(L_{i,\alpha}(x)),\qquad
  L_{i,\alpha}(D_{j\beta,i\alpha}(y))=L_{j,\beta}(y)\bmod m_i.
\]
For two distinct format classes \(\alpha,\beta\) at the same width, the same
logarithm-coordinate rule gives the automorphism
\[
  L_{i,\beta}(C_{i\alpha,i\beta}(x))=L_{i,\alpha}(x).
\]
Thus FP16 and BF16, or E4M3 and E5M2, are not silently identified even
though they share a bit width and a Field prime.

\begin{proposition}[The Field cast tower]
\label{prop:field-casts}
For \(i<j<k\), and for any choices of format classes at those widths, Field
widening maps compose, Field narrowing maps compose, and narrowing after
widening is the identity.  Suppressing the format labels in the notation,
\[
  U_{jk}U_{ij}=U_{ik},\qquad
  D_{ji}D_{kj}=D_{ki},\qquad
  D_{ji}U_{ij}=\mathrm{id}.
\]
Also \(D_{kj}U_{ik}=U_{ij}\).  In the opposite order,
\(U_{ij}D_{ji}\) is generally not the identity on \(\F_{p_j}^{\times}\); it is
the projection onto the chosen order-\(m_i\) subgroup.
\end{proposition}

\begin{proof}
The maps \(U_{ij}\) and \(D_{ji}\) are homomorphisms because they are additive in
logarithm coordinates.  The narrowing identities are just transitivity of
reduction modulo
\(m_i\mid m_j\mid m_k\).  The widening identity follows from uniqueness in the
CRT description: stepwise widening from \(i\) to \(j\) to \(k\) gives a
logarithm congruent to the original one modulo \(m_i\), and congruent to zero
modulo each step cofactor; because those cofactors are pairwise coprime, this is
the same as being zero modulo the total cofactor \(m_k/m_i\), which is the
direct widening \(U_{ik}\).  Similarly, \(D_{kj}U_{ik}\) has logarithm
\(\sigma_{ik}(a)\bmod m_j\), which is congruent to \(a\) modulo \(m_i\) and to
zero modulo \(m_j/m_i\); by uniqueness this is \(\sigma_{ij}(a)\), hence
\(D_{kj}U_{ik}=U_{ij}\).

The identity \(D_{ji}U_{ij}=\mathrm{id}\) follows because
\(\sigma_{ij}(a)\equiv a\pmod {m_i}\).  Finally, \(U_{ij}D_{ji}\) first forgets
the logarithm modulo the cofactor \(c_{ij}\) and then chooses the CRT lift with
cofactor coordinate zero, so it is a projection, not the identity in general.
The coprime-cofactor hypothesis is exactly what is used here: it gives the CRT
sections \(\sigma_{ij}\), and it turns ``zero modulo each step cofactor'' into
``zero modulo their product,'' so stepwise widening agrees with direct widening.
\end{proof}

\subsection{Actual constants}

The actual Field-family prime choices are as follows.  Variant \(1\) is used by
\texttt{Field}, \FieldFast{}, and \FieldMC{}; variant \(2\) is used by
\texttt{Field2}, \texttt{FieldFast2}, and \texttt{FieldWithMulCasts2}.  The
\(4\)-bit prime is shared because it is the only \(11\bmod 12\) prime at that
width.  The \(6\)-bit prime is standalone and is not part of the cast tower.
\[
\begin{array}{c|r@{\quad}r|r@{\quad}r|c}
  & \multicolumn{2}{c|}{\text{variant 1}}
  & \multicolumn{2}{c|}{\text{variant 2}}
  & \text{cast tower}\\
w & p & g & p & g & \\
\hline
4  & 11 & 2  & 11 & 2 & \text{yes}\\
6  & 59 & 2  & 47 & 5 & \text{standalone}\\
8  & 191 & 19 & 131 & 2 & \text{yes}\\
16 & 65171 & 2 & 64871 & 7 & \text{yes}\\
32 & 4284862331 & 2 & 4291215371 & 2 & \text{yes}\\
64 & 18446743887391934171 & 2 & 18446743217995397111 & 7 & \text{yes}
\end{array}
\]
Here \(g\) is the primitive root used for multiplicative casts in
\FieldMC{}.  At a bit width with two supported format classes, the
implementation uses \(g\) for the first class and \(g^3\) for the second:
FP8 E4M3 uses \(g\), FP8 E5M2 uses \(g^3\), FP16 uses \(g\), and
BF16 uses \(g^3\).  Since every Field-family prime is \(11\bmod 12\), \(3\) is
coprime to \(p-1\), so \(g^3\) is again a primitive root.  This is the concrete
mechanism that makes same-width casts such as FP16\(\leftrightarrow\)BF16
visible to \FieldMC{} instead of collapsing them to the identity.

\section{Why finite fields cannot model exponential laws}

In any finite field \(\F_q\), with \(q\) any prime power, any group homomorphism
\[
  (\F_q,+)\longrightarrow(\F_q^\times,\cdot)
\]
is trivial: the source has order \(q\), while the target has order \(q-1\).
Their orders are coprime.  Thus a field-valued fingerprint cannot have a
nontrivial function satisfying
\[
  \exp(a+b)=\exp(a)\exp(b)
\]
as a function only of its field residue.

To model the exponential law, the fingerprint must remember another residue
modulo a prime \(d\) for which the multiplicative side has an order-\(d\)
subgroup, alongside a prime \(p\) whose unit group contains that subgroup.  The
Chinese Remainder Theorem keeps this compact.  With \(n=pd\),
\[
  \Z/n\Z \simeq \F_p\times \F_d.
\]
The cost is algebraic: \(\Z/n\Z\) has zero-divisors, so it is a ring rather than
a field.

\section{The Sophie Germain ring}

\texttt{Semantics::SophieGermainRing} takes its target ring to be
\(\Z/n\Z\).  The modulus has the form
\[
  n=pd,\qquad p=2d+1,
\]
with \(p\) and \(d\) prime.  Strictly speaking, \(d\) is the Sophie Germain
prime and \(p\) is the corresponding safe prime; the semantics name refers to
the pair.

For widths below \(8\), the implementation does not use a composite Sophie
Germain ring.  There is too little same-width room for two sentinels and a useful
CRT exp/log channel.  Instead, the \(4\)- and \(6\)-bit cases use the
corresponding Field-family prime: the target is just \(\F_p\), the auxiliary
\(d\)-channel is absent, and exp/log operations are represented as tags.  Thus
sub-\(8\)-bit \texttt{SophieGermainRing} behaves like the faithful Field
semantics for finite algebraic arithmetic, but not for exp/log laws.

At widths \(8,16,32,64\), the actual Sophie Germain choices are as follows;
variant \(1\) is the unsuffixed semantics and variant \(2\) is
\texttt{SophieGermainRing2}.  The choices require \(d\equiv11\pmod {12}\),
equivalently \(p\equiv23\pmod {24}\).  This keeps the perfect cube-root
condition and also makes both \(2\) and \(3\) positive in the selected qr-order
factor \(p\).  At 8 bits only one pair satisfying these constraints fits, so
variant 2 intentionally shares variant 1 there.
\begin{center}
\small
\begin{tabular}{c|c|r|r|r}
variant & \(w\) & \(d\) & \(p=2d+1\) & \(n=pd\)\\
\hline
1 & 8  & 11 & 23 & 253\\
2 & 8  & 11 & 23 & 253\\
1 & 16 & 179 & 359 & 64261\\
2 & 16 & 131 & 263 & 34453\\
1 & 32 & 46199 & 92399 & 4268741401\\
2 & 32 & 45971 & 91943 & 4226711653\\
1 & 64 & 3036998843 & 6073997687 & 18446723947803676141\\
2 & 64 & 3036998183 & 6073996367 & 18446715930127601161
\end{tabular}
\end{center}

Since \(d\mid p-1\), the group \(\F_p^\times\) contains a subgroup of order
\(d\).  The implementation chooses an element \(g\in(\Z/n\Z)^\times\) whose
\(\F_p\)-component has order \(d\) and whose \(\F_d\)-component is \(1\).  It
then defines
\[
  \exp(x)=g^{x\bmod d}.
\]
This is a genuine homomorphism from the additive \(d\)-channel to the
multiplicative order-\(d\) subgroup:
\[
  \exp(a+b)=\exp(a)\exp(b).
\]

The logarithm is the dual map.  It projects a unit to the order-\(d\) subgroup in
the \(\F_p\)-component, takes the discrete logarithm to base \(g\), and returns
the corresponding element of the additive order-\(d\) subgroup of \(\Z/n\Z\).
Thus, on the modeled image,
\[
  \log(xy)=\log(x)+\log(y),\qquad
  \exp(\log(\exp x))=\exp x.
\]

The same channel carries other bases:
\[
  \exp_b(x)=g^{K_b x\bmod d}.
\]
The constants \(K_b\) are fixed units modulo \(d\).  They are not numerical
approximations to \(\log_b e\); those constants are irrational and have no
meaning as residues of dyadic values.  The contract is that each base preserves
its own law and has its own inverse.

Because the modulus is composite, division is exact only on units.  If the
denominator is a zero-divisor, the implementation returns the algebraic NaN
fingerprint.  At the sizes used here, zero-divisors are rare compared with the
zero-divisors in Triton's \(\Z/2^w\Z\), but they are real.

For order and absolute value, Sophie Germain rings use the safe-prime CRT
factor \(p=n/d\) as the selected factor in the qr-order of
Section~\ref{sec:order-abs-selection}.  Since \(p=2d+1\) with odd \(d\), this
factor is \(3\bmod4\), so the selected-factor sign convention is available at
the composite widths.  The stronger \(p\equiv23\pmod {24}\) choice above makes
\(2\) and \(3\) quadratic residues in that selected factor.  At sub-\(8\)-bit
widths, the Field-family fallback uses its prime modulus as the selected factor.

\section{The Pythagorean ring}

\texttt{Semantics::PythagoreanRing} takes its target ring to be
\(\Z/n\Z\).  The modulus has the form
\[
  n=pd,\qquad p=4d+1,
\]
again with \(p\) and \(d\) prime.  This still gives \(d\mid p-1\), so the
exp/log channel remains available.  The new feature is that
\(p\equiv 1\pmod 4\), so \(-1\) has a square root in \(\F_p\).  This allows a
finite-ring analogue of complex rotations in the \(p\)-component of the CRT
decomposition.

For widths below \(8\), the implementation does not use a composite Pythagorean
ring.  As in the Sophie Germain case, the \(4\)- and \(6\)-bit cases use the
corresponding Field-family prime, with no auxiliary \(d\)-channel.  Exp/log and
sin/cos operations are therefore represented as tags at those widths.  The usual
Pythagorean tradeoff discussed below, including the loss of perfect
\texttt{cbrt}, begins only at width \(8\); below that, finite algebraic
arithmetic and roots follow the Field-family behavior.

At widths \(8,16,32,64\), the actual Pythagorean choices are as follows; variant
\(1\) is the unsuffixed semantics and variant \(2\) is
\texttt{PythagoreanRing2}.  The checked-in choices require
\(d\equiv7\pmod {24}\).  This makes \(2\) positive in the selected qr-order
factor \(d\).  Since \(p=4d+1\) prime forces \(d\equiv1\pmod 3\) for
\(d>3\), these same choices necessarily make \(3\) a non-residue in \(d\).  At
8 bits only one pair satisfying the 2-positive constraint fits, so variant 2
intentionally shares variant 1 there.
\begin{center}
\small
\begin{tabular}{c|c|r|r|r}
variant & \(w\) & \(d\) & \(p=4d+1\) & \(n=pd\)\\
\hline
1 & 8  & 7 & 29 & 203\\
2 & 8  & 7 & 29 & 203\\
1 & 16 & 127 & 509 & 64643\\
2 & 16 & 79 & 317 & 25043\\
1 & 32 & 32647 & 130589 & 4263339083\\
2 & 32 & 32479 & 129917 & 4219574243\\
1 & 64 & 2147479879 & 8589919517 & 18446679324986898443\\
2 & 64 & 2147479447 & 8589917789 & 18446671903297182683
\end{tabular}
\end{center}

The implementation writes the construction in the quadratic algebra
\((\Z/n\Z)[i]\) with a formal \(i^2=-1\).  Here
\(\F_p\) should be read as a CRT factor of \(\Z/n\Z\), not as a unital subring:
\[
  \Z/n\Z \simeq \F_p\times \F_d.
\]
A square root of \(-1\) in the \(\F_p\)-factor gives an element whose square is
\((-1,0)\), not \((-1,-1)\).  To have a square root of \(-1\) in all of
\(\Z/n\Z\), one would also need \(-1\) to be a square modulo \(d\), which is not
part of the design constraint.  What the construction needs is weaker and more
targeted: in
\[
  (\Z/n\Z)[i]\simeq \F_p[i]\times\F_d[i],
\]
the \(\F_p[i]\) factor splits and supplies the nontrivial rotation, while the
\(\F_d[i]\) factor is kept at the identity.  The formal \(i\) keeps the
real/imaginary bookkeeping uniform across both CRT factors.  The implementation
chooses an element
\[
  \omega=\omega_{\mathrm{re}}+i\,\omega_{\mathrm{im}}
\]
of order \(d\) and norm \(1\), and defines
\[
  \cos x = \Re(\omega^{x\bmod d}),\qquad
  \sin x = \Im(\omega^{x\bmod d}).
\]
Complex multiplication immediately gives
\[
\begin{aligned}
  \cos(a+b) &= \cos a\cos b-\sin a\sin b,\\
  \sin(a+b) &= \sin a\cos b+\cos a\sin b,\\
  \cos^2 x+\sin^2 x &= 1.
\end{aligned}
\]

This variant gives up perfect cube roots.  If \(d>3\) is prime and
\(p=4d+1\) is also prime, then \(d\not\equiv 2\pmod 3\), since that would make
\(p\equiv 0\pmod 3\).  Thus \(d\equiv 1\pmod 3\), so \(3\mid d-1\).  In the
small \(d=3\) case, \(3\mid p-1\).  In either case the unit-group exponent has a
factor of \(3\), so the cube map is not globally invertible by a power map.
\texttt{cbrt} is therefore represented as a tag in the Pythagorean semantics.

For order and absolute value, Pythagorean rings use the \(d\) CRT factor as the
selected factor.  At sub-\(8\)-bit widths, the Field-family fallback uses its
prime modulus as the selected factor.  All checked-in composite-width choices
have \(d\equiv7\pmod {24}\), so the selected-factor sign convention is
available throughout the Pythagorean family and \(2\) is positive.  The same
prime-shape constraint that enables the Pythagorean construction makes \(3\) a
selected-factor non-residue in the composite cases.

\section{Tags and genuinely opaque functions}

Not every operation has a useful algebraic law in the target ring.  For
opaque libdevice calls, miscellaneous transcendentals, and Field-mode
exp/log/trig calls, the implementation produces a deterministic tag.  In this
document, a tag means the resulting fingerprint, while the operation identifier
or symbol hash is the salt used to construct it.

For a \(k\)-ary opaque operation symbol \(f\), one should think of the
implementation as choosing a deterministic map
\[
  \tau_f:R^k\longrightarrow R
\]
for the current target ring \(R\), extended to the sentinel values by the
algebraic NaN rules.  It promises congruence:
\[
  f=g,\quad a_i=b_i\ \text{for all }i
  \quad\Longrightarrow\quad
  \tau_f(a_1,\ldots,a_k)=\tau_g(b_1,\ldots,b_k).
\]
It also uses the operation identity and argument order, so changing the symbol
or changing the operands normally changes the fingerprint, up to ordinary
finite-ring collisions.  What it deliberately does not promise is an algebraic
law such as \(\log(xy)=\log x+\log y\), unless the semantics has a specific
channel implementing that law.

This is not merely a weak fallback.  Algebraically, an opaque call is best
viewed as adjoining a fresh formal generator \(t_{f,\mathbf a}\) for the
operation and its argument tuple, then continuing the surrounding calculation in
the target ring.  A fixed-width sanitizer cannot literally adjoin infinitely
many independent generators, so the tag is the compact residue fingerprint of
that formal-generator idea.  The same expression receives the same tag;
unrelated opaque expressions are kept distinct except for the unavoidable
collision risk of any finite fingerprint.

Triton uses the same opaque-operation idea for operations it does not give a
structured algebraic model.  For a fixed list of named unary operations, such as
\texttt{log}, \texttt{sqrt}, \texttt{erf}, \texttt{floor}, and \texttt{ceil},
Triton applies a hash-like permutation of the already-scrambled payload, salted
by the operation identifier.  For arbitrary extern or libdevice calls, it uses a
stable hash of the symbol name together with an argument-order-sensitive mix of
the operand payloads.  These are tags in the same semantic sense as above:
deterministic, operation-distinguishing, and deliberately uninterpreted.
\texttt{hip-fpsan} matches those formulas in \texttt{Semantics::Triton}; in the
algebraic modes the concrete mixing formula changes to live inside the target
ring, but the contract is the same.

No transcendence conjecture is involved in tags.  They behave like
uninterpreted operation symbols with congruence, not like real transcendental
functions.  No identity connecting \(\tau_f(a_1,\ldots,a_k)\) back to the
operands is asserted in the first place.

Once produced, a tag is just another fingerprint.  The ambient ring laws still
apply to expressions built from it: for example \(t+t=2t\) and
\((t+u)v=tv+uv\) are still enforced.  The only thing withheld is a mathematical
relation between \(t\) and the operands of the opaque operation.

\section{Infinity and NaN}

The finite residue model is extended by one unsigned infinity and one absorbing
NaN.  In the Field case, the finite residues together with infinity form the
projective line
\[
  \mathbb{P}^1(\F_p)=\F_p\cup\{\infty\},
\]
and the full target set of the semantics is
\[
  \mathbb{P}^1(\F_p)\cup\{\mathrm{NaN}\}.
\]
For the composite ring modes, read the same picture as the affine residue ring
together with one formal point at infinity, again plus NaN.  Correspondingly,
the arithmetic rules model the projective line rather than IEEE-754-style signed
infinities:
\[
  x/0=\infty\quad (x\ne 0),\qquad
  x/\infty=0,
\]
and indeterminate forms such as \(0/0\), \(0\cdot\infty\), and
\(\infty/\infty\) produce NaN.  Since the infinity is unsigned,
\(\infty+\infty\) also produces NaN in this model.  Arithmetic operations on
finite residues never overflow; among the basic arithmetic operations on finite
values, the only source of infinity is \(x/0\) with \(x\ne 0\).

\section{Order, absolute value, and selection}
\label{sec:order-abs-selection}

No finite ring can carry a total order compatible with its addition and
multiplication.  Consequently the algebraic semantics do not model IEEE-754
numeric ordering.  Comparisons, \texttt{min}, \texttt{max}, and \texttt{abs}
are therefore fingerprint operations, not magnitude operations.

The implementation uses the quadratic-residue split as a canonical sign
convention whenever it can select a prime factor \(q\equiv3\pmod4\).  In the
Field family, \(q\) is the whole prime modulus.  In the Sophie Germain rings,
\(q=p=n/d\), the safe-prime CRT factor.  In the Pythagorean rings, \(q=d\).
The checked-in algebraic modulus tables select such a factor in every case.
They also tune the composite factors so \(2\) is qr-positive in both Sophie
Germain and Pythagorean families.  In the Sophie Germain family, \(3\) is
qr-positive as well.  In the Pythagorean composite family, \(3\) is necessarily
a selected-factor non-residue: the condition \(p=4d+1\) with \(p,d\) prime
forces \(d\equiv1\pmod3\) for \(d>3\), while the 2-positive qr-order choice
uses \(d\equiv7\pmod {24}\).

Write \(\pi_q\) for reduction to the selected factor.  In this section,
``quadratic residue'' means a nonzero square after applying
\(\pi_q\colon\Z/n\Z\to\F_q\).  For composite rings this is deliberately not the
same as being a square in every CRT component of \(\Z/n\Z\).  Requiring all
components to be squares would fail to choose a sign in many cases, because the
unselected factor can have the opposite quadratic-residue status for \(x\) and
\(-x\).  The selected factor is a deterministic sign convention, not a full
ordered-ring structure.

Thus a finite payload \(x\) is positive when its selected component is a
nonzero quadratic residue in \(\F_q\): \(\pi_q(x)=y^2\) for some nonzero
\(y\in\F_q\).  Zero is not positive; neither are infinity, NaN, values whose
selected component is zero, or selected-factor non-residues.  Since
\(q\equiv3\pmod4\), \(-1\) is not a square in \(\F_q\), so exactly one of
\(x\) and \(-x\) is positive whenever \(\pi_q(x)\ne0\).  In the Field family
this covers every finite nonzero value.  In the composite rings, finite
zero-divisors with \(\pi_q(x)=0\) lie outside the positive class for both signs.
This gives a canonical absolute value on the selected-factor nonzero case:
\[
  |x|_{\mathrm{qr}} =
  \begin{cases}
    x, & x\text{ positive},\\
    -x, & \text{otherwise},
  \end{cases}
\]
with \(|x|_{\mathrm{qr}}=|-x|_{\mathrm{qr}}\) whenever \(\pi_q(x)\ne0\).  In
the Field-family case, in terms of the Legendre symbol
\(\chi(x)\in\{\pm1\}\), this is \(|x|_{\mathrm{qr}}=\chi(x)x\) for finite
nonzero \(x\).  The extension to the special cases is by the same
implementation rule: since zero and the sentinels are not positive,
\texttt{abs} negates them, and negation fixes them in this model.  Thus
\(\texttt{abs}(0)=0\), \(\texttt{abs}(\infty)=\infty\), and
\(\texttt{abs}(\mathrm{NaN})=\mathrm{NaN}\).

This sign convention is compatible with multiplication whenever the selected
components are nonzero.  Products of two positive residues and products of two
negative residues are positive, while a positive times a negative is negative:
\[
  Q\cdot Q\subset Q,\qquad N\cdot N\subset Q,\qquad Q\cdot N\subset N,
\]
where \(Q\) is the set of quadratic residues and \(N\) is the set of
non-residues in \(\F_q^\times\).  Equivalently,
\(\chi(\pi_q(xy))=\chi(\pi_q(x))\chi(\pi_q(y))\).  Consequently
\[
  |xy|_{\mathrm{qr}}=|x|_{\mathrm{qr}}\,|y|_{\mathrm{qr}}
\]
when \(\pi_q(x)\) and \(\pi_q(y)\) are nonzero.  Every square with nonzero
selected component is positive.  In the Field-family case, the square-root
power map \(S\) from Constraint~\ref{con:sqrt} is exact on squares in the
canonical sense:
\[
  S(x^2)=|x|_{\mathrm{qr}},\qquad |x|_{\mathrm{qr}}^2=x^2.
\]
For composite rings, the root power maps are the componentwise maps described
above; the qr-order only adds a selected-factor sign convention and does not
remove the usual zero-divisor caveat.

The associated comparison is only a two-class preorder:
\[
  a < b \quad\Longleftrightarrow\quad a\notin Q\ \text{and}\ b\in Q.
\]
Here \(Q\) still means only finite values with positive selected component;
zero, infinities, NaNs, selected-component zero values, and selected-factor
non-residues all lie outside \(Q\).  The non-strict comparison is not the same
relation as the strict comparison.  It is defined in the usual C++ way,
\(a\le b\) iff not \(b<a\), hence
\[
  a \le b \quad\Longleftrightarrow\quad b\in Q\ \text{or}\ a\notin Q.
\]
Therefore \(a<b\) implies \(a\le b\), but many pairs satisfy \(a\le b\) without
satisfying \(a<b\).  For example \(0\le n\) and \(n\le0\) both hold for a
finite selected-factor non-residue \(n\), while neither \(0<n\) nor \(n<0\)
holds.

The \texttt{min} and \texttt{max} operations use the same positive/non-positive
split, with deterministic payload tie-breaking inside a class.  This tie-break
is a selection convention, not an additional comparison law.  In particular,
\texttt{max} prefers a positive residue over zero or a non-residue, and all
finite values with nonzero selected component become positive after
\texttt{abs}.

This convention is not a magnitude order, and it does not make ordered-field
algebra possible.  It is nevertheless useful operationally in rank-revealing
linear algebra.  Pivoting code often first needs the algebraic property that a
chosen pivot or normalization denominator is nonzero; the real-valued fact that
it is the largest magnitude candidate is secondary for FPSan's exact-equivalence
purpose.  With the \(\mathrm{qr}\) absolute value,
\(\max_i |a_i|_{\mathrm{qr}}\) is nonzero
whenever any candidate has nonzero selected component.  Likewise, if a routine
compares candidate denominators represented as exact squares \(r_i^2\), the
candidates with nonzero selected component are positive, and in the Field-family
case the chosen square root is the canonical \(|r_i|_{\mathrm{qr}}\).  In
composite modes, the same selected-factor convention helps avoid artificial
order divergences, while division still has the genuine ring caveat: a selected
value that is a zero-divisor in the full ring produces the algebraic NaN
fingerprint.  This is the algebraic analogue of choosing a largest pivot or
largest \(L^2\)-norm denominator: the selection is deterministic and preserves
the square/root/nonzero relationships available in the selected factor, while
making no claim to reproduce numerical magnitude order.

Triton FPSan uses the signed-payload fallback path instead: comparisons,
\texttt{min}, and \texttt{max} use signed-payload order, and \texttt{abs}
selects between \(x\) and \(-x\) by comparing \(x\) with zero in that same
order.

\section{Extension to irrational number schemes}

The construction above used one special fact about ordinary binary floating
point: finite values are dyadic rationals, hence elements of the localization
\(\Z[1/2]\).  Logarithmic number systems are a long-standing alternative
representation \cite{swartzlander-alexopoulos,alam-garland-gregg,nguyen-solovyev-gopalakrishnan},
and many LNS-like variants make irrational algebraic values exactly
representable by allowing fractional binary exponents: half-integral exponents
give \(\sqrt2\), quarter-integral exponents give \(2^{1/4}\), and finer grids
go further.  Such formats may or may not keep a mantissa.  Multiplication is
then exponent addition, while addition is the hard operation.  For FPSan
reduction, the only essential change is that the finite values now lie in a
number field rather than in \(\Q\).

In general, let \(K\) be a number field with ring of integers \(\mathcal O_K\),
and suppose the representable finite values lie in a localization
\(S^{-1}\mathcal O_K\), or in an order inside it.  For every maximal ideal
\(\mathfrak p\) avoiding \(S\), reduction gives
\[
  S^{-1}\mathcal O_K \longrightarrow
  \mathcal O_K/\mathfrak p.
\]
The target is a finite field, not necessarily a prime field.  Thus ``reduce
floats mod \(p\)'' is shorthand for reducing the localized ring of algebraic
integers modulo a finite prime.  The algebraic number theory used here is
standard; see Neukirch \cite[Chapter I]{neukirch} for prime ideals,
localization, ramification, and the Dedekind--Kummer theorem.

Consider the concrete quarter-exponent example.  Let
\[
  \alpha = 2^{1/4},\qquad K=\Q(\alpha),\qquad \alpha^4=2.
\]
A format with scale factors \(2^{k/4}\) and a binary mantissa has finite values
in the localized order
\[
  \Z[1/2,\alpha].
\]
Negative quarter-exponents are represented using
\[
  \alpha^{-1}=\alpha^3/2,
\]
so they introduce no new primes in the denominator beyond \(2\).  The same
description covers the presence or absence of a binary mantissa, and any
mantissa width, since mantissa bits introduce only integer coefficients and
additional powers of \(2\) in the denominator.  The polynomial \(X^4-2\) is
Eisenstein at \(2\), and the discriminant of the basis
\(1,\alpha,\alpha^2,\alpha^3\) is
\[
  \operatorname{disc}(X^4-2)=4^4\,2^3=2^{11}.
\]
Thus ramification is confined to \(2\).  Every odd prime remains available for
reduction, exactly as for dyadic rationals.

For an odd rational prime \(p\), the Dedekind--Kummer theorem reads the
decomposition of \((p)\) from the factorization of \(X^4-2\) over \(\F_p\)
\cite[Chapter I]{neukirch}:
\[
\begin{array}{c|c|c}
\text{condition on }p & X^4-2\text{ over }\F_p
  & \text{residue fields}\\
\hline
p\equiv 1\pmod 8,\ 2\text{ a fourth power}
  & 1+1+1+1
  & \F_p,\F_p,\F_p,\F_p\\
p\equiv 1\pmod 8,\ 2\text{ not a fourth power}
  & 2+2
  & \F_{p^2},\F_{p^2}\\
p\equiv 7\pmod 8
  & 1+1+2
  & \F_p,\F_p,\F_{p^2}\\
p\equiv 3\pmod 8
  & 2+2
  & \F_{p^2},\F_{p^2}\\
p\equiv 5\pmod 8
  & 4
  & \F_{p^4}
\end{array}
\]
Here \(1+1+2\) means two linear factors and one irreducible quadratic factor.
The first line is the split case.  The last line is the inert case: \(p\)
remains prime in \(\Z[\alpha]\), and the residue field is the non-prime finite
field \(\F_{p^4}\).  The degree-two cases similarly produce residue fields
\(\F_{p^2}\).  Only the completely split case lands entirely in copies of
\(\F_p\).

For fingerprinting, any of these prime ideals can be a target.  If its residue
degree is \(f\), the finite target field is \(\F_{p^f}\), and the leaf map sends
\[
  m\,\alpha^k\,2^{-r}
  \longmapsto
  \overline m\,\overline\alpha^{\,k}\,\overline 2^{-r}
  \quad\text{in }\F_{p^f}.
\]
Since \(p\) is odd, \(\overline2\) is invertible, so the scale factors remain
units and the same zero-safety argument applies.  This is not implemented here,
but the algebraic mechanism is unchanged: ordinary binary floats use the number
field \(\Q\); quarter-exponent LNS uses \(\Q(2^{1/4})\); other algebraic scale
systems use other number fields.  Maximal ideals not inverted by the format
still supply finite fields in which exact algebraic identities can be
fingerprinted.

\begin{thebibliography}{9}

\bibitem{triton-fpsan-source}
Triton language and compiler.
\emph{Triton source tree}.
GitHub repository.
\url{https://github.com/triton-lang/triton}.

\bibitem{goucher-fpsan-blog}
Adam P. Goucher.
\emph{Schanuel's Conjecture and the Semantics of FPSan}.
Triton author's blog, 3 May 2026.
\url{https://cp4space.hatsya.com/2026/05/03/schanuels-conjecture-and-the-semantics-of-fpsan/}.

\bibitem{ocp-ofp8}
Open Compute Project.
\emph{OCP 8-bit Floating Point Specification (OFP8), Revision 1.0}.
1 December 2023.
\url{https://www.opencompute.org/documents/ocp-8-bit-floating-point-specification-ofp8-revision-1-0-2023-12-01-pdf-1}.

\bibitem{swartzlander-alexopoulos}
Earl E. Swartzlander, Jr. and Aristides G. Alexopoulos.
\emph{The Sign/Logarithm Number System}.
IEEE Transactions on Computers, C-24(12):1238--1242, 1975.

\bibitem{alam-garland-gregg}
Suneel A. Alam, John Garland, and David Gregg.
\emph{Low-precision logarithmic number systems: beyond base-2}.
arXiv:2102.06681, 2021.
\url{https://arxiv.org/abs/2102.06681}.

\bibitem{nguyen-solovyev-gopalakrishnan}
Tuan-Hung Nguyen, Alexey Solovyev, and Ganesh Gopalakrishnan.
\emph{Rigorous Error Analysis for Logarithmic Number Systems}.
arXiv:2401.17184, 2024.
\url{https://arxiv.org/abs/2401.17184}.

\bibitem{neukirch}
Jürgen Neukirch.
\emph{Algebraic Number Theory}.
Grundlehren der mathematischen Wissenschaften, volume 322.
Springer, 1999.

\end{thebibliography}

\end{document}
